% !TEX program = xelatex
% Шаблон научной статьи НИТУ «МИСиС»
% Заполните блоки ниже и скомпилируйте: xelatex article.tex

\documentclass{misisarticle}

% --- Классификационные коды (оставьте пустым {}, если не требуются) ---
\grnti{28.17.19}
\asjc{2604, 1706}
\oecd{1.01, 1.02}

% --- Название ---
\title{Исследование режимов самоорганизации в стохастической дискретной системе и прогнозирование её макродинамики}

% --- Авторы: \misisauthor{ФИО}{аффилиация}{город, страна}{email} ---
\misisauthor{В.~С.~Соляник}{кафедра инженерной кибернетики НИТУ «МИСиС»}{Москва, Россия}{m2305805@edu.misis.ru}
\misisauthor{А.~И.~Тихонов}{кафедра инженерной кибернетики НИТУ «МИСиС»}{Москва, Россия}{aitispu@mail.ru}

% --- Аннотация и ключевые слова ---
\abstract{%
  исследуется стохастическая дискретная система с локально взаимодействующими агентами, пространственной памятью и наследуемым поведением. Цель работы — выявить зависимость макроскопических режимов системы от параметров модели и оценить возможность прогнозирования её динамики по временным рядам. Предлагается проводить вычислительные эксперименты с независимыми случайными реализациями, анализировать энтропийные и структурные характеристики и обучить модель для прогнозирования выбранных макропоказателей. Качество прогнозирования оценивается по MAE и RMSE на отложенных данных; устойчивость выявленных закономерностей проверяется статистическим анализом повторных запусков.
}

\keywords{%
  стохастическая динамическая система, самоорганизация, фазовая диаграмма, энтропийный анализ, нейросетевое прогнозирование, статистическое моделирование
}

\begin{document}

\maketitle

% ============================================================
\section{Введение}

Проблема возникновения устойчивых структур из хаотического движения большого числа элементов остаётся одной из центральных в синергетике и неравновесной термодинамике~\cite{1,1}. Дискретные имитационные модели позволяют связать локальные правила поведения элементов с интегральными статистическими характеристиками системы и проверять такие связи вычислительным экспериментом. В качестве базовой аналогии обычно берут термодинамическую систему: поведение отдельной частицы трудно предсказать, но состояние ансамбля описывается энтропией и информационными мерами~\cite{1,1}. Развитием этого подхода служат модели искусственной жизни, в том числе непрерывные клеточные автоматы~\cite{1}, и нейроэволюция~\cite{1}, в которых поведение элементов задаётся наследуемыми параметрами.

Вычислительные работы такого рода часто опираются на единичные запуски стохастической модели. Макропоказатели подобных систем флуктуируют, поэтому визуально различимые кривые могут лежать в пределах шума, а выводы о влиянии параметров"--- оказаться артефактом случайной реализации. Требуются серии независимых запусков, доверительные интервалы и проверка гипотез с поправкой на множественные сравнения~\cite{holm}. Кроме того, для систем с нестационарной динамикой (релаксация, приток энергии, рождение и гибель структур) остаётся открытым вопрос, насколько предсказуема их макродинамика по временным рядам и дают ли нейросетевые модели~\cite{1,1} выигрыш по сравнению с тривиальными прогнозами.

Цель работы"--- выявить зависимость макроскопических режимов стохастической дискретной системы с локально взаимодействующими агентами, пространственной памятью и наследуемым поведением от параметров модели и оценить возможность прогнозирования её динамики по временным рядам. Вклад работы состоит в следующем:
\begin{itemize}
  \item вычислительное исследование зависимости режимов самоорганизации от параметров модели;
  \item статистическая оценка устойчивости результатов на независимых случайных запусках;
  \item разработка и проверка подхода к прогнозированию макродинамики системы по временным рядам на траекториях, не участвовавших в обучении.
\end{itemize}

% ============================================================
\section{Основная часть}

\subsection{Подраздел}

Подзаголовки второго уровня набираются \textit{курсивом}.

Пример маркированного списка:
\begin{itemize}
  \item первый пункт;
  \item второй пункт;
  \item третий пункт.
\end{itemize}

Пример формулы в строке: мера Жаккара $IoU = \dfrac{|A \cap B|}{|A \cup B|}$.

Нумерованная формула (ссылка в тексте: формула~\eqref{eq:f1}):
\begin{equation}
  F_1 = 2 \cdot \frac{P \cdot R}{P + R}
       = \frac{2TP}{2TP + FP + FN},
  \label{eq:f1}
\end{equation}
где $P$~--- precision, $R$~--- recall, $TP$, $FP$, $FN$~--- число истинно и ложно положительных/отрицательных детекций.

% --- Рисунок ---
\begin{figure}[htbp]
  \centering
  % Замените заглушку на свой файл:
  % \includegraphics[width=0.75\linewidth]{figures/figure}
  \fbox{\parbox[c][3cm][c]{0.9\columnwidth}{\centering Placeholder для рисунка}}
  \caption{Пример иллюстрации. Замените заглушку на свой файл из папки figures.}
  \label{fig:example}
\end{figure}

% --- Таблица ---
\begin{table}[htbp]
  \caption{Пример таблицы с результатами}
  \label{tab:example}
  \centering
  \begin{tabular}{lcc}
    \toprule
    \textbf{Метод} & \textbf{Precision} & \textbf{Recall} \\
    \toprule
    Метод~A & 0.85 & 0.78 \\
    Метод~B & 0.91 & 0.83 \\
    \bottomrule
  \end{tabular}
\end{table}

% ============================================================
\section{Заключение}

Кратко сформулируйте основные выводы работы.

% ============================================================
% Список литературы — вариант 1: вручную
\begin{thebibliography}{9}

\bibitem{example2024} Author A., Author B. Title of the paper. \textit{Journal Name}, vol.~1, no.~2, pp.~100--110, 2024.

\bibitem{book2023} Author C. \textit{Book Title}. Publisher, City, 2023.

\end{thebibliography}

\end{document}
